\subsubsection{Cosmos 3 Reasoner as Prompt Upsampler}
\label{sec::prompt_upsampling}

Prompt upsampling is the pre-generation reasoning step in Cosmos 3 that translates compact user intent into a structured spatiotemporal scene specification. Rather than merely rewriting the prompt, it expands sparse user input into a physically grounded control language that captures scene layout, temporal evolution, and audio cues when applicable. Users typically provide a short natural-language request, while high-quality image and video generation depends on many details that are rarely specified explicitly such as subject attributes, spatial layout, camera behavior, lighting, temporal ordering, audio cues, and generation controls such as resolution, aspect ratio, duration, and frame rate. The upsampler fills this gap by acting as an instruction-following LLM module between the user interface and the generator. It takes the brief request, the output's structure template, plus the optional conditioning input signals such as a starting image for image-to-video, and produces a dense typed JSON scene specification that the downstream image or video model can render from.

The prior work falls into two threads. The first thread maps short user queries into richer text prompts for image generation. DALL-E 3 uses descriptive re-captioning and caption upsampling to better match the generator's caption distribution \citep{betker2023improving}; Prompt Expansion explicitly learns to expand a query into multiple optimized text-to-image prompts \citep{datta-etal-2024-prompt}; and prompt-adaptation systems such as Promptist, BeautifulPrompt, and RePrompt train language models to rewrite user inputs into model-preferred prompts while optimizing image-level objectives \citep{hao2023optimizing,cao-etal-2023-beautifulprompt,wu2025reprompt}. The second thread uses LLMs as planners for downstream visual generators, producing object layouts, grounded scene descriptions, frame-level prompts, or multi-scene video plans \citep{lian2024llmgrounded,feng2023layoutgpt,hong2023direct2v,
lin2023videodirectorgpt}. Cosmos 3 takes the same basic idea of inserting a reasoning module before generation, but further changes the interface. The upsampler emits one typed multimodal scene program, rather than only a rewritten prompt or a generator-specific layout plan, and the program spans image, video, and audio-conditioned generation. In a sense, the upsampler proceeds by first imagining the scene, considering the temporal/spatial aspects, and adding further multimodal content (\eg, audio descriptors) in an abstract language description state, then translating them into the structured output that matches the provided template.

We treat upsampling as complex \emph{multimodal} instruction following with \emph{physical priors} over \emph{structured data}. The input may be text-only, image-plus-text, or video-plus-text. The output is a schema-constrained JSON control program that captures the semantic content, visual attributes, task-specific caption, temporal plan for video tasks, audio-description field for video tasks, and generation parameters needed by the renderer. The structured output is not only a formatting choice, it exposes the latent variables that the generator must condition on, including entities, spatial relations, actions, timing, camera behavior, audio consequences, and generation controls. This connects prompt upsampling to the body of work on structured LLM reasoning and verification, where natural-language instructions are converted into explicit, checkable programs over structured fields \citep{chegini-etal-2025-repanda}. At the same time, filling a dense scene schema requires reasoning under uncertainty in which the upsampler must infer likely object states, temporal transitions, causal interactions, and acoustic outcomes from incomplete user input \citep{pournemat2025reasoning} in order to expand along the creativity axis without contradicting physical constraints of imposed reality. In Cosmos 3, these probabilistic and physical priors are expressed through the schema itself. The model first imagines a coherent world state, then maps it into a temporal rollout, and finally derives audio cues that remain synchronized with visible events. This separation is useful because it exposes prompt understanding as an independently inspectable component rather than entangling it with the rendering model.

\paragraph{Prompt contract.}
At inference time, the upsampler receives the user description together with the selected generation controls. For image-to-video generation, it also receives the conditioning image, which is treated as definitive visual evidence for the first frame. The request is formatted as four tagged blocks: 1) \texttt{instructions} block: introduces all blocks, defines the task-level behavior such as producing a single fenced JSON object, avoiding extra commentary, and treating the template as mandatory. 2) \texttt{image} or \texttt{video} \texttt{\_description} block: carries the raw semantic request; for conditioned requests such as I2V, the attached condition content is provided alongside this text and the constraints specify how visual facts should be anchored to it. 3) \texttt{task\_constraints} block: carries most of the
operational detail. It specifies field ordering, timing format, duration bounds, copied controls such as resolution, aspect ratio, duration, and frame rate, preservation requirements for user-provided entities and actions, and task-specific grounding rules such as matching the I2V first frame at \(t=0\). 4) \texttt{output\_json\_template} block: defines the schema surface: which keys must be populated, which nested fields are expected, and where scene, temporal, audio, subject, camera, lighting, style, and control information should be placed. In other words, the constraints describe how the upsampler should reason and copy values, while the JSON template defines the shape that the final answer
must take. The system prompt is a minimal ``You are a helpful assistant''. The full template can be found in appendix~\ref{appendix:upsampler_prompt_template}. Using this template, certain placeholder arguments such as description, FPS, duration, aspect ratio, resolution, \etc, have to be filled accordingly before passing to the Upsampler.
